\chapter{Método do trabalho} \label{cap:metodo}

O desenvolvimento é incremental e organiza-se em seis fases. Elas não são
estanques, já que a revisão bibliográfica prossegue durante a implementação e a
especificação é revisitada sempre que uma medição contraria uma suposição, mas
cada uma delas produz um resultado que as fases seguintes consomem. Os
resultados aparecem nos Capítulos~\ref{cap:especificacao} e
\ref{cap:desenvolvimento}, e não neste capítulo.

\section{Revisão da literatura}
\label{sec:fase-revisao}

A primeira fase produz o Capítulo~\ref{cap:conceitos}. Além de situar o
trabalho, ela cumpre uma função prática. Como três das quatro ferramentas de
referência não estão disponíveis, quase todo número publicado precisa ser
tratado como afirmação do autor, e a revisão serve para separar o que pode ser
reimplementado daquilo que só pode ser citado.

\section{Estabelecimento do gabarito}
\label{sec:fase-gabarito}

Antes de qualquer análise, as anotações do Racebench são lidas, contadas e
certificadas contra uma contagem manual, e tanto a unidade de contagem quanto a
regra que decide se um candidato reportado corresponde a um defeito anotado são
fixadas por escrito. Medir recall exige um denominador, e a
Seção~\ref{sec:benchmarks-conceito} mostrou que a literatura utiliza
denominadores diferentes sem justificar a escolha. Esta fase precede a
implementação precisamente para que nenhuma decisão de contagem possa ser
tomada após a observação dos resultados.

\section{Especificação}
\label{sec:fase-especificacao}

A terceira fase produz os requisitos do Capítulo~\ref{cap:especificacao},
acompanhados da lista de premissas sob as quais a ferramenta afirma não perder
defeitos. A lista é mantida junto ao código e revisada a cada alteração que
afete a análise.

\section{Implementação da etapa estática}
\label{sec:fase-estatica}

A quarta fase implementa a etapa estática em três estágios de custo crescente,
na ordem em que executam. Cada estágio só é considerado concluído quando todos
os defeitos anotados permanecem presentes em sua saída. Esse teste, aqui
denominado portão de recall, é executado a cada alteração e trata uma perda de
recall como defeito de implementação, e não como parâmetro a ajustar.

\section{Implementação da etapa com modelo de linguagem}
\label{sec:fase-llm}

A quinta fase é desenvolvida sobre um conjunto de registros de contexto
montados manualmente a partir do Racebench, descrito na
Seção~\ref{sec:conjunto-triagem}, e não sobre a saída da etapa estática. Essa
escolha permite que as duas etapas avancem em paralelo e dá ao conjunto um
segundo uso, o de servir como sonda para verificar se um \llm\ é capaz de
julgar esse tipo de defeito antes que se invista na construção completa da
etapa. A integração entre as duas ocorre ao final desta fase, e os números
medidos sobre o conjunto montado manualmente não são transportados para a
população real, pelas razões discutidas na Seção~\ref{sec:ameacas}.

\section{Avaliação}
\label{sec:fase-avaliacao}

A sexta fase segue o protocolo da Seção~\ref{sec:avaliacao}, aplicado primeiro
ao Racebench e depois à suíte de 18 programas reais.
